\chapter{CHAPTER 3: METHODOLOGY}\label{chapter-3-methodology}

\section{3.1 Chapter Overview}\label{chapter-overview}

This chapter delineates the research methodology and development
framework employed to design, develop, and evaluate the SentinelAQ
predictive system. The methodology is structurally aligned with
Saunders' Research Onion to provide a rigorous philosophical and
practical foundation for the study. Section 3.2 unpacks the theoretical
research framework, justifying the philosophical stance and strategic
approach. Section 3.3 outlines the Agile development methodology
utilized for system implementation. Section 3.4 details the project
management lifecycle, including a comprehensive Gantt chart, phase
deliverables, and resource requirements. Section 3.5 presents a critical
risk analysis and mitigation strategy. Finally, Section 3.6 details the
quantitative evaluation design, specifying the metrics used to benchmark
model accuracy and interpretability.

\section{3.2 Research Methodology Framework (Saunders' Research
Onion)}\label{research-methodology-framework-saunders-research-onion}

To ensure academic rigor and coherence, this research systematically
addresses the layers of Saunders' Research Onion (Saunders et al.,
2019), tailoring each methodological choice to specifically address the
research aim of forecasting and interpreting PM2.5 variations in Sri
Lanka.

\subsection{3.2.1 Research Philosophy:
Positivism}\label{research-philosophy-positivism}

\textbf{Choice \& Justification:} This study adopts a
\textbf{Positivist} philosophy. Positivism relies on observable,
quantifiable, and empirical data to generate statistically verifiable
outcomes. Because this research aims to forecast atmospheric pollutant
concentrations using mathematical algorithms (XGBoost, LSTM, SARIMAX)
applied to objective sensor telemetry (PM2.5, temperature, humidity), it
operates under the assumption of an objective, measurable reality. This
stance perfectly aligns with Research Objective 4 (evaluating model
accuracy via RMSE/MAE), as it requires a strictly quantitative,
value-free interpretation of data rather than subjective human
interpretation.

\subsection{3.2.2 Research Approach:
Deductive}\label{research-approach-deductive}

\textbf{Choice \& Justification:} A \textbf{Deductive} approach is
utilized. Deductive research begins with existing theories and
formulates hypotheses to be tested empirically. Based on the literature
review (Chapter 2), the study deduced the hypothesis that tree-based
ensemble models (XGBoost) configured with tabular lag features would
outperform sequential deep learning models (LSTM) in short-to-medium
horizon environmental forecasting. The research design is structured
explicitly to test this deduction through rigorous algorithmic
benchmarking across 1 to 48-hour horizons.

\subsection{3.2.3 Research Strategy: Design Science Research (DSR) \&
Experimentation}\label{research-strategy-design-science-research-dsr-experimentation}

\textbf{Choice \& Justification:} The strategy is a hybrid of
\textbf{Design Science Research (DSR)} and \textbf{Experimentation}. DSR
focuses on creating novel, practical artifacts to solve specific
problems---in this case, the design and development of the SentinelAQ
automated pipeline and mobile application (Objective 7). Concurrently,
an Experimental strategy is employed within the machine learning phase
(Objectives 3 \& 4), where variables are controlled (e.g., standardizing
the 48-feature vector across all models) to benchmark competing
algorithms against a statistical baseline (SARIMAX).

\subsection{3.2.4 Time Horizon:
Longitudinal}\label{time-horizon-longitudinal}

\textbf{Choice \& Justification:} A \textbf{Longitudinal} time horizon
is required. Air quality is inherently cyclical, heavily influenced by
seasonal monsoons, daily traffic patterns, and annual temperature
shifts. To capture these non-linear temporal dynamics, the dataset must
span a continuous, multi-year period. Consequently, data was collected
longitudinally from January 2022 to December 2025, allowing the models
to learn long-term seasonal dependencies and short-term diurnal spikes.

\subsection{3.2.5 Data Collection and Analysis
Techniques}\label{data-collection-and-analysis-techniques}

\textbf{Choice \& Justification:} The research employs
\textbf{Quantitative Data Collection and Analysis}. It must be
explicitly noted that the primary research methodology for this study is
entirely empirical; no qualitative user interviews or surveys were
conducted. Instead, primary data is programmatically collected via
Application Programming Interfaces (APIs)---specifically the PurpleAir
REST API for historical PM2.5 ground truth and the Open-Meteo API for
localized weather proxies. The analysis technique involves supervised
machine learning regression, statistically evaluated using error
metrics, and mathematically interpreted using cooperative game theory
(SHAP).

\section{3.3 System Development
Methodology}\label{system-development-methodology}

To manage the complexities of data engineering, algorithmic tuning, and
full-stack software development, an \textbf{Agile (Iterative
Prototyping)} methodology was adopted.

\textbf{Justification:} Traditional Waterfall methodologies are too
rigid for machine learning projects, where data quality issues or model
underperformance often necessitate revisiting earlier stages. Agile
allows for iterative cycles (Sprints). For instance, when initial deep
learning models failed to capture 48-hour trends accurately, the Agile
framework permitted a ``step back'' to the feature engineering phase to
introduce 24-hour rolling statistical means, rapidly testing the new
vector on XGBoost in the subsequent sprint.

\subsection{3.3.1 Development Sprints}\label{development-sprints}

The project was divided into five core iterative phases: 1.
\textbf{Phase 1 (Data Acquisition \& Preprocessing):} Automating API
extraction, handling missing values, and chronologically aligning
PurpleAir (dependent variable) and Open-Meteo (independent variables)
datasets for Colombo and Kandy. 2. \textbf{Phase 2 (Feature Engineering
\& Baseline):} Constructing the 48-dimensional feature vector (lags,
rolling means, cyclical encoding) and establishing the statistical
forecasting baseline using SARIMAX. 3. \textbf{Phase 3 (Advanced
Modeling \& Optimization):} Iteratively training, tuning, and
cross-validating the XGBoost and LSTM architectures. 4. \textbf{Phase 4
(Explainability \& Inference Pipeline):} Implementing the SHAP
TreeExplainer for the optimal model and wrapping the inference logic
into an automated Python backend connected to Firebase. 5. \textbf{Phase
5 (Application Development):} Building the React Native mobile front-end
to visualize the Firebase payload (AQI dials, SHAP insights, push
notifications).

\section{3.4 Project Management and
Timeline}\label{project-management-and-timeline}

\subsection{3.4.1 Gantt Chart}\label{gantt-chart}

The following Mermaid diagram outlines the project timeline, showcasing
major phases, milestones, and dependencies.

\begin{Shaded}
\begin{Highlighting}[]
\NormalTok{gantt}
\NormalTok{    title SentinelAQ Project Timeline}
\NormalTok{    dateFormat  YYYY{-}MM{-}DD}
\NormalTok{    axisFormat  \%b \%Y}
    
\NormalTok{    section Foundation}
\NormalTok{    Project Proposal \& Literature Review :done, p1, 2025{-}01{-}15, 30d}
\NormalTok{    Mid{-}Point Review Submission         :done, p2, 2025{-}02{-}15, 14d}
    
\NormalTok{    section Phase 1 \& 2: Data \& Features}
\NormalTok{    API Integration \& Data Extraction   :done, d1, 2025{-}03{-}01, 14d}
\NormalTok{    Data Cleaning \& Preprocessing       :done, d2, after d1, 10d}
\NormalTok{    Feature Engineering (48{-}vector)     :done, d3, after d2, 14d}
\NormalTok{    SARIMAX Baseline Training           :done, d4, after d3, 7d}
    
\NormalTok{    section Phase 3: ML Modeling}
\NormalTok{    XGBoost Training \& Tuning           :done, m1, after d3, 14d}
\NormalTok{    LSTM Architecture \& Training        :done, m2, after m1, 14d}
\NormalTok{    Model Benchmarking \& Selection      :done, m3, after m2, 7d}
    
\NormalTok{    section Phase 4: Pipeline \& XAI}
\NormalTok{    SHAP Implementation                 :active, x1, after m3, 10d}
\NormalTok{    Firebase Backend Integration        :active, x2, after x1, 10d}
    
\NormalTok{    section Phase 5: Mobile App}
\NormalTok{    React Native UI Design              :active, a1, after x2, 14d}
\NormalTok{    App Integration \& Testing           :a2, after a1, 14d}
    
\NormalTok{    section Finalization}
\NormalTok{    Thesis Writing (Chapters 4{-}7)       :w1, 2025{-}06{-}01, 30d}
\NormalTok{    Final Presentation Prep             :w2, after w1, 10d}
\end{Highlighting}
\end{Shaded}

\subsection{3.4.2 Project Deliverables}\label{project-deliverables}

\begin{itemize}
\tightlist
\item
  \textbf{Phase 1/2:} Cleaned, synchronized CSV datasets for Colombo and
  Kandy; Engineered 48-feature training matrices.
\item
  \textbf{Phase 3:} Trained model artifacts (\texttt{.pkl},
  \texttt{.h5}); CSV files containing comprehensive RMSE/MAE metrics
  across all horizons.
\item
  \textbf{Phase 4:} Automated Python inference script; Active Firebase
  Firestore database populated with real-time predictions; SHAP summary
  visualizations.
\item
  \textbf{Phase 5:} Functioning React Native mobile application
  (SentinelAQ) deployed on an emulator/device.
\item
  \textbf{Final:} Bound Thesis Report and Presentation Slide Deck.
\end{itemize}

\subsection{3.4.3 Resource Requirements}\label{resource-requirements}

\begin{itemize}
\tightlist
\item
  \textbf{Hardware:} A local development workstation with sufficient RAM
  (16GB+) and a multi-core CPU for parallel processing of XGBoost trees
  and LSTM matrix multiplications.
\item
  \textbf{Software/Languages:} Python 3.11+ (Data Science Backend),
  JavaScript/React Native (Mobile Frontend).
\item
  \textbf{Libraries:} \texttt{scikit-learn}, \texttt{xgboost},
  \texttt{tensorflow}/\texttt{keras}, \texttt{statsmodels},
  \texttt{shap}, \texttt{pandas}, \texttt{numpy}, \texttt{matplotlib}.
\item
  \textbf{APIs \& Cloud:} PurpleAir REST API, Open-Meteo API, Google
  Firebase (Firestore \& Cloud Messaging).
\end{itemize}

\section{3.5 Risk Management}\label{risk-management}

Proactive risk management was essential to ensure project completion.
The following table identifies technical and methodological risks, their
potential impact, and the mitigation strategies executed.

{\def\LTcaptype{none} % do not increment counter
\begin{longtable}[]{@{}
  >{\raggedright\arraybackslash}p{(\linewidth - 6\tabcolsep) * \real{0.2500}}
  >{\raggedright\arraybackslash}p{(\linewidth - 6\tabcolsep) * \real{0.2500}}
  >{\raggedright\arraybackslash}p{(\linewidth - 6\tabcolsep) * \real{0.2500}}
  >{\raggedright\arraybackslash}p{(\linewidth - 6\tabcolsep) * \real{0.2500}}@{}}
\toprule\noalign{}
\begin{minipage}[b]{\linewidth}\raggedright
Risk Identification
\end{minipage} & \begin{minipage}[b]{\linewidth}\raggedright
Likelihood
\end{minipage} & \begin{minipage}[b]{\linewidth}\raggedright
Impact
\end{minipage} & \begin{minipage}[b]{\linewidth}\raggedright
Mitigation Strategy
\end{minipage} \\
\midrule\noalign{}
\endhead
\bottomrule\noalign{}
\endlastfoot
\textbf{R1: Satellite Cloud Occlusion} & High & Critical &
\textbf{Mitigated:} Pivoted methodology from Sentinel-5P satellite
fusion to robust ground-based meteorological proxies (Open-Meteo) to
ensure zero data gaps in the inference pipeline. \\
\textbf{R2: API Rate Limits \& Failures} & Medium & High &
\textbf{Mitigated:} Implemented local CSV caching during the training
phase. The live inference pipeline includes \texttt{try-except} blocks
to handle transient API timeouts gracefully. \\
\textbf{R3: Overfitting on Sparse Data} & High & High &
\textbf{Mitigated:} Utilized strict chronological Train/Test splits (no
random shuffling, which prevents time-series data leakage). Applied
regularization parameters in XGBoost and Dropout layers in LSTM. \\
\textbf{R4: SHAP Computational Bottleneck} & Medium & Medium &
\textbf{Mitigated:} Selected \texttt{TreeExplainer} for the XGBoost
model, which is mathematically optimized for tree structures and
computes Shapley values exponentially faster than generic Kernel
explainers. \\
\end{longtable}
}

\section{3.6 Evaluation Design}\label{evaluation-design}

A rigorous evaluation framework is necessary to validate whether the
research objectives were achieved and the research questions answered.

\subsection{3.6.1 Validation Methods}\label{validation-methods}

To evaluate the predictive accuracy of the models (Objective 4), the
dataset was partitioned using a strict chronological split. Data from
2022 to 2024 was allocated to the training set, while the complete
calendar year of 2025 was sequestered as an unseen test set. Random
cross-validation (like K-Fold) was explicitly avoided, as shuffling
time-series data causes ``look-ahead bias'' (data leakage), artificially
inflating model accuracy.

\subsection{3.6.2 Evaluation Metrics}\label{evaluation-metrics}

The models are mathematically benchmarked using two primary loss
functions and one practical classification metric: 1. \textbf{Root Mean
Squared Error (RMSE):} The primary metric for environmental forecasting.
Because RMSE squares the differences before averaging, it heavily
penalizes large errors. In a public health context, failing to predict a
massive toxic smog spike is significantly more dangerous than multiple
minor miscalculations; minimizing RMSE ensures the model is structurally
tuned to detect severe anomalies. 2. \textbf{Mean Absolute Error (MAE):}
Provides the linear, average magnitude of forecasting errors. It offers
an intuitive metric (e.g., ``The model is off by an average of 4.3
µg/m³''), making baseline performance easily interpretable for
non-technical stakeholders. 3. \textbf{R-Squared (R²) and AQI Category
Accuracy:} R² measures the proportion of variance explained by the
model. Furthermore, numeric predictions are converted into standardized
US EPA Air Quality Index (AQI) categories (e.g., ``Good'',
``Unhealthy'') to calculate a percentage-based categorical accuracy,
reflecting real-world utility.

\subsection{3.6.3 Explainability
Evaluation}\label{explainability-evaluation}

The Explainable AI component (Objective 5) is evaluated using SHAP
Summary Plots and Dependence Plots. The evaluation criterion is not a
numerical error metric, but rather logical cohesion: do the top SHAP
features align with atmospheric physics? For example, validating whether
the model successfully identifies that high wind speeds decrease PM2.5
(dispersion) while historical PM2.5 lags exponentially increase the
likelihood of future spikes (accumulation).
